% Generated draft: requires mathematical and source review.
\documentclass[11pt]{article}
\usepackage[margin=25mm]{geometry}
\usepackage{fontspec}
\setmainfont{Latin Modern Roman}
\usepackage{amsmath,amssymb,mathtools,unicode-math}
\setmathfont{Latin Modern Math}
\usepackage{enumitem,needspace,xurl,hyperref}
\setlength{\parindent}{0pt}
\setlength{\parskip}{5pt}
\setlength{\emergencystretch}{3em}
\newcommand{\N}{\mathbb{N}}
\newcommand{\Z}{\mathbb{Z}}
\newcommand{\Q}{\mathbb{Q}}
\newcommand{\R}{\mathbb{R}}
\newcommand{\C}{\mathbb{C}}
\newcommand{\F}{\mathbb{F}}
\newcommand{\A}{\mathbb{A}}
\newcommand{\PP}{\mathbb{P}}
\newcommand{\E}{\mathbb{E}}
\newcommand{\eps}{\varepsilon}
\newcommand{\dd}{\,\mathrm{d}}
\begin{document}
\section*{UnsolvedMath ID: 701024}
\textbf{Source problem code:} catalogue entry
\subsection{Definitions and mathematical statement}
Fix a prime \(p\). A pro-\(p\) group is a topological group isomorphic to an inverse limit of finite \(p\)-groups. Let \(V\) be a variety of pro-\(p\) groups, and let \(F(V)\) be a free pro-\(p\) group in \(V\), as in the supplied source. The source does not specify the free rank or its convention for varieties of pro-\(p\) groups and their free objects.

An element \(g\in F(V)\) is called periodic if it has finite order. Every such order is a power of \(p\). The question is whether the orders of the periodic elements of \(F(V)\) are bounded. Equivalently, does there exist an integer \(n\geq 0\) such that
\[
\forall g\in F(V),\qquad
\bigl(\exists m\geq 1\colon g^m=1\bigr)
\ \Longrightarrow\ g^{p^n}=1?
\]
Here the bound is sought for the given \(F(V)\); the source does not explicitly say whether it must be independent of the free rank. No bound uniform over different varieties or primes is explicitly requested.

The supplied source states that the answer is affirmative when \(V\) is metabelian, meaning that every group in \(V\) has abelian commutator subgroup. It attributes this special case to A. N. Zubkov's 1997 doctoral dissertation. This special-case assertion is retained as a source report and has not been independently verified.
\subsection{Short English statement}
Must the finite-order elements of a free group in a variety of pro-\(p\) groups have uniformly bounded orders? The source reports an affirmative answer for metabelian varieties. The free rank and whether the bound must be independent of that rank are unspecified.
\subsection{Sources}
\begin{itemize}
\item Local UnsolvedMath registry, permanent ID \(701024\), exact source problem code \texttt{catalogue entry}, title \emph{Kourovka Notebook Problem 14.39}. The supplied record labels the problem \emph{open}; this is a supplied classification, not an independently verified assessment.
\item \emph{The Kourovka Notebook (Unsolved Problems in Group Theory)}, Problem 14.39, proposed by A. N. Zubkov in the 14th issue (1999). The supplied background also identifies the 21st issue edition (2026), \texttt{arXiv:1401.0300v48}. The supplied collection description names E. I. Khukhro and V. D. Mazurov as editors. These bibliographic details have not been independently checked.
\item A. N. Zubkov, \emph{Dr. Sci. Diss.}, Omsk, 1997 (Russian), cited by the supplied statement for the metabelian special case. No dissertation title, page reference, or link is supplied.
\end{itemize}
\end{document}
